%%%%%
%
%% Dans ce fichier, les �ventuels accents sont cod�s
%% en ISO-8859-15 (aka latin9)
%% La ligne suivante donne trois exemples 
%% e accent aigu : �, e accent grave : �, e-dans-l'o : �
\section{Vent 3D � partir du vent apparent en coordonn�es rectangulaires et de la vitesse du support dans le rep�re terrestre}

\itemstyle{Op\'erateur~:} \operfmt{CalcVent\_rect\_repgeo}
\index{CalcVent\_rect\_repgeo}

%%%%%OBJET%%%%
\itemstyle{Objet~:}
Cet op\'erateur calcule les 3 composantes du vecteur vent
dans le rep�re g�ographique terrestre, � partir de mesures sur
une plate-forme mobile, dans le cas o� :
\begin{itemize}
 \item la mesure du vent relatif � la plate-forme ("vent apparent")
est exprim�e en coordonn�es rectangulaires.
 \item la mesure de la vitesse de la plate-forme par rapport au sol
est exprim�e dans le rep�re g�ographique local (vers l'Est, le Nord, et 
la verticale).
\end{itemize}

%%%%%ARGUMENTS%%%%
\bigskip
%% Dans ce fichier, les �ventuels accents sont cod�s
%% en ISO-8859-15 (aka latin9)
%% La ligne suivante donne trois exemples 
%% e accent aigu : �, e accent grave : �, e-dans-l'o : �
%%%%%ENTREES%%%%
\begin{tabarg}
\multicolumn{3}{l}{\titretab {Entr\'ees}}\\
\hline
\$Arg1 & Coeff., \texttt{float} & blabla (m)\\
\$Arg2 & Coeff., \texttt{float} & blabla (m)\\
\$Arg3 & Coeff., \texttt{float} & blabla (m)\\
\$Arg4 & Vecteur, \texttt{float} & blabla (m/s)\\
\$Arg5 & Vecteur, \texttt{float} & blabla (m/s)\\
\$Arg6 & Vecteur, \texttt{float} & blabla (m/s)\\
\$Arg7 & Vecteur, \texttt{float} & blabla (m/s)\\
\$Arg8 & Vecteur, \texttt{float} & blabla (m/s)\\
\$Arg9 & Vecteur, \texttt{float} & blabla (m/s)\\
\$Arg10 & Vecteur, \texttt{float} & blabla (radian/s)\\
\$Arg11 & Vecteur, \texttt{float} & blabla (radian/s)\\
\$Arg12 & Vecteur, \texttt{float} & blabla (radian/s)\\
\$Arg13 & Vecteur, \texttt{float} & blabla (radian)\\
\$Arg14 & Vecteur, \texttt{float} & blabla (radian)\\
\$Arg15 & Vecteur, \texttt{float} & blabla (radian)\\
\end{tabarg}

\medskip
%%%%%SORTIES%%%%
\begin{tabarg}
\multicolumn{3}{l}{\titretab {Sorties}}\\
\hline
\$Arg16 & Vecteur, \texttt{float} & blabla (\$Arg4.unit)\\
\$Arg17 & Vecteur, \texttt{float} & blabla (\$Arg4.unit)\\
\$Arg18 & Vecteur, \texttt{float} & blabla (\$Arg4.unit)\\
\end{tabarg}


%%%%%FORMULE%%%%
\itemstyle{Formule}

%% Formule de l'operateur : CalcVent_rect_repgeo
%% Dans ce fichier, les �ventuels accents sont cod�s
%% en ISO-8859-15 (aka latin9)
%% La ligne suivante donne trois exemples 
%% e accent aigu : �, e accent grave : �, e-dans-l'o : �
\[
Resultat = a-b. \Zeta
\]



%%%%%AUTEUR%%%%
\itemstyle{Auteur~:}
CNRM/GMEI/TRAMM

%%%%%REFERENCE%%%%
\itemstyle{R\'ef\'erences~:}
Bulletin NCAR-RAF Nr 23 \cite{RAF-23}.
